 \section*{第 \ref{ch:b2:rigid-body-mechanics} 章——刚体力学}

\begin{solution}{exo:b2:rigid-body-mechanics:1}
(a) $\tfrac1{12}M\ell^2 = \qty{4.2e-2}{kg.m^2}$；绕一端（惠更斯，
$d = \ell/2$）：$\tfrac13M\ell^2 = \qty{0.167}{kg.m^2}$。(b) $\tfrac12MR^2 =
\qty{4.0e-2}{kg.m^2}$；绕边缘：$\tfrac32MR^2 = \qty{0.12}{kg.m^2}$。
(c) $\tfrac25 \times 0.16 \times (0.0285)^2 = \qty{5.2e-5}{kg.m^2}$。(d) 圆环，
$MR^2$ 与 $\tfrac12MR^2$ 相比：两倍。
\end{solution}

\begin{solution}{exo:b2:rigid-body-mechanics:2}
$v = \qty{5.0}{m/s}$，$\omega = v/R = 5/0.35 = \qty{14.3}{rad/s}$。顶部：
$2v = \qty{10}{m/s}$，向前；底部：$0$；前部：$\vect v_G + \vect\Omega
\wedge\vect{GM} = (v, -v)$，即 $v\sqrt2 = \qty{7.1}{m/s}$，方向与水平成
\qty{45}{\degree} 向下。相对自行车，轮胎各处的速度均为
$R\omega = v = \qty{5}{m/s}$，处处如此。石子以 \qty{10}{m/s}
的速度水平向前飞出。
\end{solution}

\begin{solution}{exo:b2:rigid-body-mechanics:3}
$J = \tfrac12 \times 40 \times 0.25^2 = \qty{1.25}{kg.m^2}$; $\omega = \qty{628}{rad/s}$.
(a) $E = \tfrac12J\omega^2 = \qty{246}{kJ}$。(b) $\Gamma = J\omega/\Delta t =
1.25 \times 628/120 = \qty{6.5}{N.m}$；末时刻功率 $P = \Gamma\omega = \qty{4.1}{kW}$。
(c) 取出 \qty{200}{kJ} 后剩余 \qty{46}{kJ}：$\omega = \sqrt{2 \times 46\,000/1.25}
= \qty{271}{rad/s} = \qty{2600}{rpm}$。(d) 转子边缘的空气阻力随
$\omega^2$ 增大（功率随 $\omega^3$ 增大），会使转子发热并减速；在
真空中、采用磁悬浮轴承时，损耗降至几瓦。
\end{solution}

\begin{solution}{exo:b2:rigid-body-mechanics:4}
$J = \tfrac12MR^2 + MR^2 = \tfrac32MR^2$，$d = R$：$\ell_{\text{等效}} = \tfrac32R
= \qty{0.30}{m}$，$T = 2\pi\sqrt{0.30/9.81} = \qty{1.10}{s}$。在 $R/2$ 处：$J =
\tfrac12MR^2 + \tfrac14MR^2 = \tfrac34MR^2$，$d = R/2$，$\ell_{\text{等效}} =
\tfrac32R$ 同样成立，周期相同。一般地，$\ell_{\text{等效}}(d) = (R^2/2 +
d^2)/d$，在 $d = R/\sqrt2$ 时取最小值：$\ell_{\text{等效}} = R\sqrt2 = \qty{0.283}{m}$，
$T = \qty{1.07}{s}$.
\end{solution}

\begin{solution}{exo:b2:rigid-body-mechanics:5}
(a) $\tan20^\circ = 0.36 \le f(1 + k)/k$，即 $\le 0.6$（圆环），$0.9$
（圆盘），$1.05$（球）：均能滚动。(b) $v = \sqrt{2gh/(1 + k)}$：圆环
\qty{4.4}{m/s}，圆盘 \qty{5.1}{m/s}，球 \qty{5.3}{m/s}；球最先到达，圆环
最后。(c) 斜面长度 $h/\sin\alpha = \qty{5.85}{m}$，$a = g\sin\alpha/(1 + k) =
1.68$, $2.24$, $\qty{2.40}{m/s^2}$, $t = \sqrt{2d/a} = 2.6$, $2.3$, \qty{2.2}{s}.
(d) $T = \dfrac{k}{1 + k}Mg\sin\alpha = \tfrac27Mg\sin\alpha$；接触
点静止，因此 $\vect T\cdot\vect v_I = 0$。(e) 圆环：$\tan\alpha >
2f = 0.6$，$\alpha > \qty{31}{\degree}$；球：$\tan\alpha > 3.5f = 1.05$，
$\alpha > \qty{46}{\degree}$.
\end{solution}

\begin{solution}{exo:b2:rigid-body-mechanics:6}
$\Gamma = 2fNr = 2 \times 0.4 \times 800 \times 0.12 = \qty{77}{N.m}$; $\dot\omega =
\Gamma/J = \qty{64}{rad/s^2}$; $\omega_0 = \qty{94}{rad/s}$, $t = \qty{1.5}{s}$;
$\theta = \omega_0^2/2\dot\omega = \qty{69}{rad} = 11$ 圈；生热 $= \tfrac12J
\omega_0^2 = \qty{5.3}{kJ}$, $\Delta T = 5300/(1.5 \times 470) = \qty{7.6}{K}$.
\end{solution}

\begin{solution}{exo:b2:rigid-body-mechanics:7}
(a) 绳子必须施加净力矩才能使滑轮加速转动：$T_1 \ne T_2$。
(b) $m_1a = m_1g - T_1$，$m_2a = T_2 - m_2g$，$J\dot\omega = (T_1 - T_2)R$，其中
$a = R\dot\omega$: $a = (m_1 - m_2)g/(m_1 + m_2 + J/R^2) = 0.5 \times 9.81/4.0 =
\qty{1.23}{m/s^2}$。(c) $T_1 = m_1(g - a) = \qty{17.2}{N}$，$T_2 = m_2(g + a) =
\qty{16.6}{N}$。(d) 滑轮无质量时：\qty{1.40}{m/s^2}。能量法：$(m_1 - m_2)gx =
\tfrac12(m_1 + m_2 + J/R^2)v^2$；对时间求导。
\end{solution}

\begin{solution}{exo:b2:rigid-body-mechanics:8}
(a) 重力 $mg$ 作用于中点；墙面：顶端水平方向的 $N_w$；地面：
竖直方向的 $N = mg$，以及水平方向的摩擦力 $T = N_w$。对
脚端取矩：$mg\,\tfrac{\ell}2\cos\theta = N_w\ell\sin\theta$。(b) $T = N_w =
mg/(2\tan\theta) \le fmg \iff \tan\theta \ge 1/2f$。(c) $\tan\theta \ge
1.25$，$\theta \ge \qty{51}{\degree}$。(d) 人位于 $x$ 处：$N = 4mg$，$N_w =
mg(\tfrac12 + 3x/\ell)/\tan60^\circ \le 4fmg = 1.6mg$: $x/\ell \le (2.77 -
0.5)/3 = 0.76$.
\end{solution}

\begin{solution}{exo:b2:rigid-body-mechanics:9}
(a) 接触点以 $v + R\omega$ 向前运动（$\omega > 0$ 对应
逆旋）：摩擦力 $-fmg$ 作用于质心，$\dot v = -fg$；力矩
$-fmgR$ 作用于自旋：$\dot\omega = -5fg/2R$。(b) 滑动速度 $v + R\omega = v_0 +
R\omega_0 - \tfrac72fgt$ 在 $t_1 = 2(v_0 + R\omega_0)/7fg$ 时为零，此后 $v_1
= v_0 - fgt_1 = (5v_0 - 2R\omega_0)/7$。(c) $v_1 < 0 \iff R\omega_0 > \tfrac52v_0$。
(d) $t_1 = 16/(7 \times 1.96) = \qty{1.2}{s}$；$v_1 = (10 - 12)/7 = \qty{-0.29}{m/s}$：
球折返。
\end{solution}

\begin{solution}{exo:b2:rigid-body-mechanics:10}
(a) 取 $x$ 指向拉力方向，摩擦力 $T$ 沿 $x$ 方向：$Ma = F + T$；
对 $G$ 取矩（逆时针为正）：$rF + RT = -kMRa$（滚动，$\omega_z
= -a/R$）。因此 $a = F(R - r)/(1 + k)MR > 0$：线轴
\emph{朝向}拉力方向滚动，将线绕上。(b) $T = -F(kR + r)/(1 +
k)R$（向后）；若 $F \le fMg(1 + k)R/(kR + r)$，则不滑动。(c) 线的
作用线与轮毂相切；其对接触
点 $I$ 的力矩为 $F(r - R\cos\beta)$，顺时针（向前滚动），当 $\cos\beta >
r/R$，作用线过 $I$ 时力矩为零——即瞬时轴
——此时 $a = F(R\cos\beta - r)/(1 + k)MR$，$T = -F(kR\cos\beta + r)/(1 + k)R$，
$N = Mg - F\sin\beta$。(d) 当 $\cos\beta < r/R$ 时，对 $I$ 的力矩为
逆时针：线轴\emph{背离}拉绳者滚动。
\end{solution}

\begin{solution}{exo:b2:rigid-body-mechanics:11}
(a) $Mv_0 = P$；对 $G$ 取矩：$\tfrac25MR^2\omega_0 = P(h - R)$（上旋方向
$h > R$ 时）。(b) $v_0 = R\omega_0 \iff 2R = 5(h - R) \iff h = \tfrac75R$；高于此值时
$R\omega_0 > v_0$（上旋），低于此值时为逆旋或下旋。(c)
$\omega_0 = 0$：滑动持续 $t_1 = 2v_0/7fg = \qty{0.44}{s}$，滑动距离 $v_0t_1 -
\tfrac12fgt_1^2 = \qty{1.1}{m}$。(d) 位于 $\tfrac75R$ 处的缓冲垫使球
无滑动滚动：没有滑动阶段，不会在台面上损失速率，
回弹可预测。
\end{solution}

\begin{solution}{exo:b2:rigid-body-mechanics:12}
(a) $v_G = (R - r)\dot\theta$；接触点静止，因此
以 $\omega = v_G/r = (R - r)\dot\theta/r$ 绕其轴自转。
(b) $E_k = \tfrac12Mv_G^2 + \tfrac12\cdot\tfrac12Mr^2\omega^2 = \tfrac34M(R - r)^2
\dot\theta^2$，$E_p = -Mg(R - r)\cos\theta$；由 $\dd E/\dd t = 0$ 可得 $\ddot\theta
+ \dfrac{2g}{3(R - r)}\sin\theta = 0$，周期 $T = 2\pi\sqrt{3(R - r)/2g}$。(c) 滑动
质点的周期为 $2\pi\sqrt{(R - r)/g}$——滚动时周期长 $\sqrt{3/2}$ 倍，
其中三分之一的能量是转动能。(d) 切向：$M(R - r)\ddot\theta =
-Mg\sin\theta + T$，于是 $T = \tfrac13Mg\sin\theta$；在转折点处
$N = Mg\cos\theta_0$：$f \ge \tfrac13\tan\theta_0 \approx \theta_0/3$。
\end{solution}

\begin{solution}{pb:b2:rigid-body-mechanics:1}
\textbf{1.} $v = \qty{25}{m/s}$，$\omega = \qty{80.6}{rad/s}$；轮顶 \qty{50}{m/s}，
轮底 $0$，前缘 $(25, -25)$：\qty{35}{m/s}，以 \qty{45}{\degree} 角斜向下。

\textbf{2.} $\tfrac12Mv^2 = \qty{375}{kJ}$；$4 \times \tfrac12J\omega^2 = \qty{13}{kJ}$；
总计 \qty{388}{kJ}，车轮占 $3.3\%$。

\textbf{3.} $N_f + N_r = Mg = \qty{11.8}{kN}$，$N_fa = N_rb$：$N_f = Mgb/L =
\qty{6.8}{kN}$，$N_r = Mga/L = \qty{5.0}{kN}$——前轴即发动机所在轴。

\textbf{4.} $P_r = \sum\mu_rN_iR\omega = \mu_rMgv = 0.012 \times 11772 \times 25 =
\qty{3.5}{kW}$；阻力 $\tfrac12 \times 1.2 \times 0.7 \times 625 = \qty{262}{N}$，\qty{6.6}{kW}。

\textbf{5.} $(M + 4J/R^2)\dot v = -\mu_rMg - \tfrac12\rho C_xSv^2$，$M_{\text{有效}} =
1200 + 41.6 = \qty{1242}{kg}$。仅滚动阻力时：$\dot v = \qty{-0.114}{m/s^2}$，
$d = (25^2 - 12.5^2)/0.228 = \qty{2.1}{km}$。仅阻力时：$v = v_0\eu^{-x/\lambda}$，
$\lambda = M_{\text{有效}}/\tfrac12\rho C_xS = \qty{2.96}{km}$，$d = \lambda\ln2 =
\qty{2.0}{km}$（两者同时作用：$\qty{1.0}{km}$）。

\textbf{6.} 重力、法向力，以及路面作用于轮胎的摩擦力。
前轮所受向前的摩擦力是唯一的水平外力，
因此它使汽车加速；其功率为零（接触点
静止）；动能来自发动机——内力做功。

\textbf{7.} $M\dot v = T_f - T_r$；前轮 $2J\dot v/R = \Gamma - T_fR$；
后轮 $2J\dot v/R = T_rR$。相加得：$(M + 4J/R^2)\dot v = \Gamma/R$；
$M_{\text{有效}} = \qty{1242}{kg}$。

\textbf{8.} 对 $G$ 取矩（水平地面力位于深度 $h$ 处，
合力为 $M\dot v$）：$aN_f - bN_r + hM\dot v = 0$，结合 $N_f + N_r = Mg$，于是
$N_r = Mga/L + Mh\dot v/L$：后轮增加 $Mh\dot v/L$，前轮则减少同样的量。

\textbf{9.} 前轮驱动：$M\dot v \le fN_f = f(Mgb/L - Mh\dot v/L)$，于是
$\dot v_{\max} = fgb/(L + fh) = 11.77/3.04 = \qty{3.9}{m/s^2}$。后轮驱动：
$M\dot v \le f(Mga/L + Mh\dot v/L)$，$\dot v_{\max} = fga/(L - fh) = \qty{4.0}{m/s^2}$。
载荷转移有利于后轴：后轮驱动时，后部的重力
使汽车的加速更为强劲。

\textbf{10.} $\Gamma = RM_{\text{有效}}\dot v_{\max} = 0.31 \times 1242 \times 3.87 =
\qty{1.5}{kN.m}$；$P = \Gamma v/R = 1490 \times 13.9/0.31 = \qty{67}{kW}$。

\textbf{11.} $\dot v_{\max} = 0.1 \times 9.81 \times 1.5/2.655 = \qty{0.55}{m/s^2}$；
需 \qty{25}{s} 才能达到 \qty{50}{km/h}。

\textbf{12.} 抓地时为静摩擦力，无滑动，接触处无耗散。
空转时为动摩擦力（摩擦因数略小），
其做功使轮胎发热，并磨光或熔化冰面，发动机的
功率用于使车轮空转，而非推动汽车。

\textbf{13.} 制动力矩为内力（车轮–车身）；汽车
仅由路面作用于轮胎的向后摩擦力减速。对车轮而言
（$J\dot\omega$ 可忽略）：$\Gamma_b \approx T_iR \le fN_iR$。

\textbf{14.} 同样的平衡，取 $\dot v = -\gamma$：$N_f = Mgb/L + Mh\gamma/L$；
当 $\gamma = fg = \qty{7.85}{m/s^2}$ 时：$N_f = 6792 + 1992 = \qty{8.8}{kN}$，$N_r =
\qty{3.0}{kN}$。

\textbf{15.} $\gamma_{\max} = fg = \qty{7.85}{m/s^2}$；$d = v^2/2\gamma = \qty{40}{m}$；
湿路面 $f = 0.45$：$\qty{4.4}{m/s^2}$，\qty{71}{m}。

\textbf{16.} 力矩相等则作用力相等，但后轴只承受
\qty{3.0}{kN}，因此先抱死。滑动轮胎的摩擦力
方向与其滑动速度相反，不能再提供侧向力：抱死的
后轴失去方向稳定性，汽车甩尾。

\textbf{17.} 抱死时：$\gamma = 0.7g = \qty{6.9}{m/s^2}$，减小 $12\%$，停车
距离为 \qty{45}{m}（而非 \qty{40}{m}），且无法转向。
防抱死系统感知车轮向抱死方向减速，松开并
每秒多次重新施加制动，使轮胎保持在
静摩擦峰值附近并滚动。

\textbf{18.} \qty{388}{kJ} 注入热容为 $4 \times 6 \times 470 = \qty{11.3}{kJ/K}$ 的制动器：$\Delta T =
\qty{34}{K}$，每次制动。一辆 \qty{40}{t} 的卡车下坡 \qty{1000}{m} 时须耗散
$Mgh = \qty{390}{MJ}$：制动器过热并衰退；砂石避险坡道使
卡车借助砂石的摩擦力停下。

\textbf{19.} $0.7 \times 388 = \qty{272}{kJ} = \qty{0.075}{kWh}$；一块 \qty{50}{kWh}
电池约可供 $660$ 次此类制动。

\textbf{20.} $Mv^2/\rho \le fMg$：$v \le \sqrt{fg\rho} = \qty{19.8}{m/s} =
\qty{71}{km/h}$；湿路面为 \qty{53}{km/h}。

\textbf{21.} 侧向摩擦力 $Mv^2/\rho$ 作用于深度 $h$ 处，法向力位于
$\pm w/2$：$(N_{\text{外}} - N_{\text{内}})\,w/2 = Mv^2h/\rho$，于是 $\Delta N =
Mv^2h/\rho w$。内侧车轮在 $\Delta N = Mg/2$ 时离地：$v = \sqrt{g\rho w/2h} =
\qty{25.9}{m/s} = \qty{93}{km/h} > \qty{71}{km/h}$：车辆先侧滑（由于 $f <
w/2h = 1.36$）。

\textbf{22.} 在倾角为 $\beta$ 的弯道上，若无摩擦，$N\sin\beta =
Mv^2/\rho$，且 $N\cos\beta = Mg$：$v = \sqrt{g\rho\tan\beta} = \qty{9.3}{m/s} =
\qty{33}{km/h}$，对应 \qty{10}{\degree} 倾角。

\textbf{23.} $\omega_0 = \qty{6.28}{rad/s}$，$J\dot\omega = \qty{-0.05}{N.m}$：$t =
\qty{126}{s}$；损失的能量 $\tfrac12J\omega_0^2 = \qty{20}{J}$。

\textbf{24.} 动摩擦力 $fN = \qty{2.4}{kN}$ 推动汽车向前：
$\dot v = 2400/1200 = \qty{2.0}{m/s^2}$；对车轮有 $J\dot\omega = -fNR$，
$\dot\omega = \qty{-744}{rad/s^2}$。滑动速度 $R\omega - v$ 初值为 \qty{1.95}{m/s}，
以 $\qty{233}{m/s^2}$ 的速率下降，经 \qty{8.4}{ms} 消失，此时汽车已
获得 \qty{1.7}{cm/s} 的速度增量，并移动了 \qty{70}{\micro m}——这是一次冲击，而非推动。

\textbf{25.} 滚动时：摩擦力 $\approx 0$，滚动阻力 $\mu_rN$，
$|T| \le fN$ 显然成立。加速时，驱动轮所受的
向前摩擦力 $T \le fN_{\text{驱动}}$，$\dot v \le fgb/(L + fh)$。制动时：
所有轮胎上的向后摩擦力，$T_i \le fN_i$，$\gamma \le fg$。
转弯时：侧向摩擦力 $Mv^2/\rho \le fMg$。各种情况下，
由一个力——路面摩擦力，受库仑摩擦定律约束——
承担全部工作。
\end{solution}
